\section{Conclusion}
This controlled extension of the fixed-relevance design of \textcite{chen_huang_tien_2021} reduces excluded-instrument signal while independent compensating noise preserves the joint treatment/outcome-disturbance distribution under the stated Gaussian construction. Comparing three preserved high-dimensional IVQR implementations shows why point accuracy, direct-test acceptance, false-candidate rejection and finite-grid informativeness must be assessed together.

The central empirical disagreement is between nominal-looking truth acceptance and useful exclusion of false effects. Retaining most of a grid while rejecting few false candidates offers little discrimination. Direct-at-truth acceptance is not membership coverage of the stored randomized grid set, and strong-relevance undercoverage remains a counterpoint to near-nominal rates elsewhere. Implementation-level RMSE advantages have documented exceptions and do not isolate regularization.

These findings establish neither implemented orthogonality nor validity of the nominal chi-square cutoff or automatic transfer of theoretical weak-identification robustness. The implementation and selected-case numerical limitations summarized in the Discussion restrict that link. Their certificates do not establish invariance of historical decisions or a full-experiment error frequency. Historical results are retained, finite grids establish no unrestricted confidence-set geometry, and no particular correction is shown to restore calibration.
